
\section{Introducción}
En las últimas décadas, el sector agrícola se enfrenta a la necesidad de optimizar el consumo de 
recursos hídricos y fertilizantes mediante el despliegue de las denominadas técnicas de agricultura de 
precisión (\textit{smart farming}). La monitorización de las condiciones fisicoquímicas del terreno resulta 
determinante para prevenir la degradación de los cultivos y minimizar el impacto medioambiental. Tradicionalmente, 
la caracterización del suelo se ha realizado mediante la recolección manual de muestras y su posterior procesado en 
laboratorios especializados, un método costoso y con tiempos de respuesta incompatibles con la toma de decisiones ágil en campo.

El paradigma del Internet de las Cosas (IoT), sustentado por microcontroladores de alta eficiencia 
y redes de comunicación de baja potencia y largo alcance (LPWAN), permite superar estas barreras. 
Este Trabajo Fin de Grado aborda el diseño, la implementación y la validación de un sistema integral de adquisición 
y análisis de parámetros fisicoquímicos del suelo. La solución se fundamenta en un nodo sensor autónomo basado en el 
módulo LilyGO T-SIM7080G provisto del microcontrolador ESP32-S3. Esta placa integra un módem SIM7080G empleado como receptor 
GNSS para la sincronización temporal periódica de la hora del sistema, corrigiendo la desviación intrínseca (\textit{drift}) 
del reloj en tiempo real (RTC) del microcontrolador sin incurrir en el elevado coste energético de mantener activo dicho módulo de forma continua.

Para maximizar la vida útil de la celda de alimentación y asegurar una operatividad 
desatendida a largo plazo, el nodo aprovecha la interfaz de recarga solar y el soporte 
de batería integrados en la placa LilyGO T-SIM7080G, gobernados por la unidad de gestión 
de potencia (PMU AXP2101). A nivel de software, el firmware opera bajo un ciclo de vida 
periódico regido por la función \texttt{setup()}, prescindiendo del bucle continuo 
\texttt{loop()} para alternar ráfagas breves de procesamiento con periodos prolongados 
de suspensión profunda (\textit{Deep Sleep}). El intervalo previsto de suspensión es de 600\,s (10 minutos), al que se suma el tiempo activo del ciclo. En cada despertar, el PMU conmuta de
forma selectiva los carriles de alimentación de los subsistemas periféricos. El nodo 
evalúa inicialmente el porcentaje de batería para abortar el ciclo si se sitúa bajo un
umbral crítico de protección; a continuación, excita el bus industrial RS-485 para 
adquirir las variables de la sonda multivariable bajo el protocolo Modbus RTU e inserta 
la muestra empaquetada (\texttt{struct Data} de 19 bytes) en un búfer circular retenido 
en la memoria RTC Slow RAM mediante el atributo \texttt{RTC\_DATA\_ATTR}, preservando 
las lecturas durante los ciclos de suspensión profunda, pero no ante la pérdida completa de alimentación del dominio RTC.

Con el propósito de optimizar el compromiso entre el consumo energético de la radio 
y las restricciones de tiempo en el aire (\textit{airtime}) impuestas por la normativa 
europea en la banda EU868, el dispositivo ejecuta un algoritmo de cálculo dinámico 
del tamaño de carga útil. Dicha lógica evalúa la capacidad máxima admitida (\textit{Max Payload}) 
en función del factor de dispersión (\textit{Spreading Factor}) y ancho de banda vigentes, 
determinando si el volumen de lecturas retenidas en la memoria RTC Slow RAM justifica 
la activación del transceptor LoRa SX1276 para proceder a su transmisión agrupada 
por lotes (\textit{batching}) en una única trama LoRaWAN activada por OTAA. 

Asimismo, la arquitectura contempla persistencia local en la tarjeta microSD (bus SD-MMC).
El protocolo de envío agrupa las muestras del búfer RTC en conjuntos pendientes y limita a dos los
intentos de cada conjunto. Si falla el primero, su identidad, tamaño e intentos consumidos se conservan
en \texttt{pending\_batch} durante la suspensión profunda. Si en el segundo intento cabe menos, se
intenta enviar el prefijo que quepa y el excedente se conserva para un lote nuevo, con dos intentos
propios. El prefijo se registra y descarta tanto si la transmisión local termina sin error como si agota
sus dos intentos. Si el segundo intento no puede enviar ninguna muestra por fallo de radio, sesión o
capacidad nula, se registra y descarta el conjunto pendiente completo. El éxito local no confirma la
recepción por la pasarela. Durante el flujo normal entre ciclos, cada conjunto finalizado se registra una
sola vez; no existe garantía de atomicidad entre SD y RTC ante un corte durante la finalización. Una
escritura parcial se detecta comparando los bytes informados con la longitud solicitada, pero puede dejar
una última línea JSON incompleta. La interfaz usada no permite truncar ese fragmento atómicamente y el
lector debe rechazar líneas incompletas. Si falla SD, el conjunto finalizado se descarta igualmente,
aceptando la posible pérdida o truncamiento del registro. La SD actúa como histórico, no como cola de
reenvío. El estado RTC distingue el conjunto
pendiente de las muestras nuevas. La retención del búfer solo se mantiene mientras el dominio RTC esté
alimentado.

Por último, la infraestructura de recepción se despliega sobre una Raspberry Pi equipada con un concentrador 
LoRa multicanal comunicado por SPI. Mediante la orquestación de servicios desatendidos administrados por el gestor 
de arranque \texttt{systemd}, la plataforma levanta de forma automática el \textit{packet forwarder}, el servidor de 
red ChirpStack, el intermediario de mensajería MQTT (Mosquitto) y un servicio de ingesta desarrollado en Python. Este 
último procesa los eventos entrantes, extrae las cargas útiles individuales o por lotes procesadas por el decodificador 
JavaScript del perfil del dispositivo e inserta los registros en un motor relacional PostgreSQL provisto de índices específicos 
sobre el identificador del nodo y la marca temporal, sirviendo como origen de datos optimizado para la consulta y análisis en tiempo 
real a través de cuadros de mando interactivos en Grafana.

\section{Motivación}
La motivación de este proyecto responde a los desafíos operativos, energéticos y de fiabilidad identificados durante la digitalización del sector agroalimentario:

\begin{itemize}
    \item \textbf{Limitaciones de las plataformas comerciales:} Gran parte de las soluciones de monitorización agrícola disponibles en el mercado presentan costes elevados, ecosistemas cerrados y dependencia de cuotas de suscripción continuadas en plataformas en la nube. Existe la necesidad de desarrollar arquitecturas abiertas e interoperables, fundamentadas en tecnologías de código abierto que faciliten la soberanía de los datos y su adaptación a múltiples tipologías de cultivo.

    \item \textbf{Autonomía energética en entornos rurales aislados:} Las explotaciones agrícolas suelen carecer de tendido eléctrico y cobertura celular regular. Garantizar una operatividad continua durante meses o años mediante baterías de capacidad moderada exige no solo implementar ciclos de suspensión profunda (\textit{Deep Sleep}), sino gestionar activamente los carriles de alimentación de los sensores y radios mediante una unidad de gestión de potencia (PMU), aislando los periféricos para eliminar corrientes parásitas o de fuga.

    \item \textbf{Optimización del canal radioeléctrico y normativa LPWAN:} El envío indiscriminado de lecturas individuales compromete la batería y colisiona con las severas restricciones de ciclo de trabajo (\textit{duty cycle}) y tiempo en el aire (\textit{airtime}) de la banda europea EU868. El diseño exige implementar estrategias de transmisión agrupada por lotes (\textit{batching}) con cálculo dinámico de carga útil, maximizando el rendimiento del espectro radioeléctrico y minimizando los encendidos del transceptor.

    \item \textbf{Tolerancia a fallos y resiliencia temporal:} Si no se obtiene sincronización GNSS válida, el nodo continúa usando el reloj del ESP32, sin garantizar que marque UTC real. Ante fallos de radio, cada conjunto pendiente dispone de dos intentos. Si la capacidad disminuye en el segundo, se finaliza solo el prefijo intentado y el excedente pasa a un lote nuevo. La SD es un histórico, no una cola de reenvío; un fallo de escritura se notifica y se acepta la posible pérdida o truncamiento del registro.

    \item \textbf{Operación desatendida en la infraestructura de recepción:} La pasarela receptora desplegada en campo debe operar como un dispositivo autónomo (\textit{appliance}), capaz de levantar la pila de comunicaciones (servidor de red, intermediario MQTT, pasarela LoRa) y la persistencia en base de datos de manera automática tras cortes imprevistos en el suministro eléctrico, sin requerir supervisión por consola.

    \item \textbf{Interés formativo y transversalidad técnica:} El proyecto ofrece la oportunidad de sintetizar y validar competencias de la Ingeniería Informática, integrando desde la programación de sistemas embebidos en C++ con restricciones en tiempo real y buses industriales (Modbus RTU sobre RS-485), hasta la ingeniería de datos en PostgreSQL, administración de servicios del sistema bajo Linux y la visualización analítica en tiempo real.
\end{itemize}

\section{Objetivos}
El \textbf{objetivo general} de este Trabajo Fin de Grado consiste en diseñar, implementar y validar un sistema integral y autónomo de monitorización edafológica basado en un nodo sensorial IoT de ultra bajo consumo (ESP32-S3) y conectividad LoRaWAN, complementado por una estación receptora sobre Raspberry Pi encargada de la ingesta desatendida, persistencia relacional y visualización web en tiempo real de los parámetros fisicoquímicos del suelo.

Para la consecución de esta meta global, se definen los siguientes \textbf{objetivos específicos}:
\begin{itemize}
    \item \textbf{Acondicionamiento y sensado industrial:} Integrar una sonda multivariable de suelo acoplada a un transceptor RS-485, implementando una máquina de consulta bajo el protocolo Modbus RTU para la adquisición de humedad, temperatura, pH, conductividad eléctrica y macronutrientes (nitrógeno, fósforo y potasio).

    \item \textbf{Gestión energética y ciclo de vida del firmware:} Desarrollar el software embebido estructurado en ráfagas de ejecución puntual (\texttt{setup()}) seguidas de suspensión profunda (\textit{Deep Sleep}), controlando selectivamente el corte y encendido de los carriles de alimentación de los periféricos mediante el PMU AXP2101 para reducir las corrientes de fuga en reposo.

    \item \textbf{Sincronización y consistencia temporal:} Integrar el módem satelital GNSS mediante comandos AT para la sincronización horaria de alta precisión en tiempo absoluto UTC, mitigando la desviación intrínseca (\textit{drift}) del reloj interno del microcontrolador sin incurrir en consumo continuo.

    \item \textbf{Optimización de radio LoRaWAN y transmisión por lotes:} Implementar la pila LoRaWAN sobre la banda EU868 con activación OTAA, incorporando un algoritmo de cálculo dinámico del tamaño máximo de carga útil para agrupar lecturas en lotes (\textit{batching}) en una estructura binaria compacta de 19 bytes, y preservando el estado de la sesión de red en memoria RTC para evitar renegociaciones tras cada suspensión.

    \item \textbf{Gestión acotada del búfer y tolerancia a fallos:} Conservar en RTC un contador de cabecera, tamaño e intentos de cada conjunto pendiente. Tras el primer fallo se conserva; tras el éxito local o el segundo intento se registra y descarta solo el conjunto finalizado. Si en el segundo intento solo cabe un prefijo, el excedente pasa a un lote nuevo sin registro ni descarte prematuros. Si no puede enviarse ninguna muestra, se finaliza el conjunto completo. Los errores SD se notifican y no impiden liberar el conjunto finalizado. La SD no es cola de reenvío; la unicidad se limita al flujo normal y no es atómica ante un corte entre escritura y actualización de RTC.

    \item \textbf{Despliegue del backend receptor y persistencia de series temporales:} Orquestar en la Raspberry Pi una pasarela LoRaWAN con concentrador SPI, integrando el servidor ChirpStack, el intermediario MQTT (Mosquitto) y un servicio puente en Python para decodificar las tramas e indexar las mediciones en una base de datos relacional PostgreSQL.

    \item \textbf{Visualización y automatización de la infraestructura:} Automatizar el arranque tolerante a fallos de todos los servicios del sistema mediante el gestor \texttt{systemd}, y diseñar cuadros de mando interactivos en Grafana con consultas SQL optimizadas para el seguimiento analítico de los cultivos.
\end{itemize}

\section{Organización de la memoria}
La presente memoria se estructura en los siguientes capítulos:
\begin{itemize}
    \item \textbf{Capítulo 2 (Estado del arte):} Examina las tecnologías de comunicación LPWAN (LoRaWAN, NB-IoT y Sigfox), analiza los protocolos de buses industriales en entornos agrícolas y revisa las principales soluciones comerciales disponibles.
    \item \textbf{Capítulo 3 (Requisitos, especificaciones, coste, riesgos y viabilidad):} Detalla los requisitos funcionales y no funcionales del sistema, el presupuesto económico estimado, la matriz de riesgos y la viabilidad técnica y operativa.
    \item \textbf{Capítulo 4 (Diseño):} Describe la arquitectura y el diseño previsto de los subsistemas.
    \item \textbf{Capítulo 5 (Implementación y pruebas):} Expone la implementación y las pruebas efectivamente realizadas, distinguiendo lo implementado de lo diseñado o pendiente.
    \item \textbf{Capítulo 6 (Conclusiones):} Resume los resultados respaldados por la implementación y las pruebas, las limitaciones y el trabajo futuro.
\end{itemize}